%summary of whole research. Include:

% Objective
% Methods
% A rough description of outcome


\section*{Abstract}

The contrast avoidance model (CAM) posits that anxiety patients attempt to avoid large shifts in emotional states and therefore maintain a heightened state of negative affect (NA). Previous research showed that this can be empirically supported using VADER sentiment estimations of Twitter data and aggregating the sentiment scores for each user, comparing an anxiety cohort with a depression cohort. In order to build upon that, the present study constructed new aggregation mechanisms. To give less weight to short-term fluctuations in NA levels, tweets are aggregated in a 6-hour rolling window and a daily window for each user separately. Then, the cohorts are compared with regards to their NA levels and variability. It was found, that depression patients show higher NA levels as well as variability when examining VADER scores. A novel analysis method based on linear mixed modeling (LMM) was proposed. LMMs were fit to the different aggregation levels and the estimated regression coefficient for the anxiety cohort was used as a measure of differences in general NA levels, while the residual standard deviation was used as a measure of unpredictability and thus NA variability that cannot be explained by available external factors. The analysis showed that this approach also supports the previous findings when applied to VADER sentiment scores, indicating that depression patients have higher NA levels and higher NA variability. When applying the proposed methods to sentiment scores obtained from transformer-based NLP tools (specifically, a RoBERTa model), the results changed. NA levels were still significantly higher for depression patients, but no differences or even opposite significant differences in NA variability were found. In the LMM approach with 6-hour windowed data, the anxiety cohort showed significantly larger NA variability than the depression cohort. Thus, the previous findings regarding empirical validation of the CAM could only be supported for VADER-based analysis and require future research in order to explain the conflicting results.
